\section{Minimum Time to Burn the Tree}
\label{sec:trees-burn}

\problemstatement{A fire starts at a given target node. Each second, the
fire spreads from every burning node to its adjacent nodes (children and
parent). Return the minimum number of seconds needed to burn the entire
tree.}

\begin{patternbox}
This is the ``distance in all directions'' idea taken to its conclusion:
the answer is the \textbf{maximum distance} from the target to \emph{any}
node, because the last node to catch fire is the farthest one. Same
recipe -- \textbf{parent map plus BFS from the target} -- but instead of
stopping at layer $K$, count how many layers it takes to exhaust the tree.
\end{patternbox}

\subsection*{Understanding the problem carefully}

Fire spreads one edge per second in every direction. A node at graph
distance $d$ from the source catches fire at second $d$. The whole tree is
ablaze once the farthest node burns, so the time equals the eccentricity
of the target: the greatest distance to any node. A BFS from the target,
counting completed layers until the queue empties, measures exactly this.

\begin{ideabox}[Answer = Number of BFS Layers to Empty the Tree]
Build the parent map, BFS from the target with a visited set. Each fully
processed frontier is one second of spread. The number of layers after the
\emph{initial} one (the source itself, at time $0$) is the burn time --
i.e.\ increment a timer each time you expand a non-empty next frontier.
\end{ideabox}

\subsection*{Algorithm}

\begin{enumerate}
  \item DFS to fill \textit{parent}[node].
  \item BFS from target; \textit{time} $=-1$ (so the first layer, the
        source, corresponds to time $0$).
  \item Each layer: mark all unvisited neighbours; if the layer expanded
        anything, increment \textit{time}. Actually, increment
        \textit{time} once per processed frontier.
  \item Return \textit{time}.
\end{enumerate}

(Concretely: initialise \textit{time} $=0$, and increment it only when a
frontier produces at least one newly burning node.)

\subsection*{Dry run}

\dryrun{Tree: root $1$; left $2$ (left child $4$); right $3$. Start fire at
$4$.}

\begin{itemize}
  \item $t=0$: $\{4\}$ burning.
  \item $t=1$: neighbour of $4$ is $2$. $\{2\}$.
  \item $t=2$: neighbours of $2$ are $1$ (and $4$ burnt). $\{1\}$.
  \item $t=3$: neighbour of $1$ is $3$. $\{3\}$.
  \item Tree fully burnt at $t=3$.
\end{itemize}

\subsection*{C++ implementation}

\begin{lstlisting}
#include <bits/stdc++.h>
using namespace std;

struct TreeNode {
    int val;
    TreeNode* left;
    TreeNode* right;
    TreeNode(int v) : val(v), left(nullptr), right(nullptr) {}
};

void markParents(TreeNode* node, unordered_map<TreeNode*, TreeNode*>& par) {
    if (node == nullptr) return;
    if (node->left)  { par[node->left]  = node; markParents(node->left,  par); }
    if (node->right) { par[node->right] = node; markParents(node->right, par); }
}

int timeToBurn(TreeNode* root, TreeNode* target) {
    unordered_map<TreeNode*, TreeNode*> par;
    markParents(root, par);

    unordered_set<TreeNode*> visited;
    queue<TreeNode*> q;
    q.push(target);
    visited.insert(target);

    int time = 0;
    while (!q.empty()) {
        int s = q.size();
        bool spread = false;
        for (int i = 0; i < s; i++) {
            TreeNode* node = q.front(); q.pop();
            for (TreeNode* nb : {node->left, node->right, par.count(node) ? par[node] : nullptr}) {
                if (nb && !visited.count(nb)) {
                    visited.insert(nb);
                    q.push(nb);
                    spread = true;
                }
            }
        }
        if (spread) time++;
    }
    return time;
}

int main() {
    TreeNode* root = new TreeNode(1);
    root->left = new TreeNode(2);
    root->right = new TreeNode(3);
    root->left->left = new TreeNode(4);

    cout << "Burn time: " << timeToBurn(root, root->left->left) << "\n";
    return 0;
}
\end{lstlisting}

\begin{complexitybox}
\textbf{Time Complexity.} $O(n)$ -- one DFS plus one BFS, each visiting
every node once. \textbf{Space Complexity.} $O(n)$ for the parent map,
visited set, and queue.
\end{complexitybox}

\begin{takeawaybox}
Burning the tree is ``nodes at distance $K$'' with $K$ pushed to the
maximum: the answer is the number of BFS layers to consume the whole tree.
Once you see fire-spread problems as level-counting BFS on the
parent-augmented graph, they become routine.
\end{takeawaybox}
